% =====================================================================
%  Distribution-free offsetting result: drop-in text for Section 4 and
%  Online Appendix B. Notation follows the manuscript (S, h, C, A, lambda,
%  s, epsilon, k, theta). Requires amsmath, amssymb, amsthm.
% =====================================================================

% ---------------------------------------------------------------------
%  Main text (replaces Corollary 1(iii) and the motivation of Prop. 2)
% ---------------------------------------------------------------------

\paragraph{Heterogeneous detectability.}
Let each defect $j$ carry its own detection rate $\rho_j>0$, drawn independently
from a distribution $G$ on $(0,\infty)$, and let a defect with rate $\rho_j$ be found
after an exponentially distributed amount of effort with that rate. Then
\[
  S(e)=\int_0^\infty e^{-\rho e}\,dG(\rho).
\]
No parametric form is imposed on $G$. For $e\ge 0$ define the
\emph{survivor distribution} $G_e$ by
$dG_e(\rho)=e^{-\rho e}dG(\rho)/S(e)$: the distribution of detection rates
among defects still undetected after effort $e$. Write $\mu_e$, $\sigma_e^2$ and
$\mathrm{CV}_e=\sigma_e/\mu_e$ for its mean, variance and coefficient of variation.
All three are finite for every $e>0$, whether or not $G$ has finite moments.

\begin{proposition}[Offsetting is set by heterogeneity in detectability]
\label{prop:cv}
Under the standing assumptions and heterogeneous detectability:
\begin{enumerate}
\item[(a)] The hazard of detection is the mean detection rate among survivors, and
  its slope is minus their variance: $h(e)=\mu_e$ and $h'(e)=-\sigma_e^2$. Hence
  $h$ is non-increasing, and strictly decreasing unless $G$ is degenerate.
\item[(b)] With $g(e)=C''(e)/C'(e)$, the elasticity in Proposition~1 is
  \[
    \varepsilon=\frac{g(e^*)+\sigma_{e^*}^2/\mu_{e^*}}
                     {g(e^*)+\sigma_{e^*}^2/\mu_{e^*}+\mu_{e^*}} .
  \]
  With linear verification cost,
  \[
    \varepsilon=\frac{\mathrm{CV}_{e^*}^2}{1+\mathrm{CV}_{e^*}^2},
    \qquad
    1-\varepsilon=\frac{1}{1+\mathrm{CV}_{e^*}^2}.
  \]
  Offsetting is therefore complete ($\varepsilon=0$) if and only if all defects are
  equally detectable. It is partial for every non-degenerate $G$, and convex cost
  only makes it more partial, because $\varepsilon$ increases in $g$.
\item[(c)] (Hardest-to-find defects govern the limit.) Suppose
  $G(\rho)=\rho^{k}L(\rho)$ as $\rho\downarrow 0$, for some $k>0$ and a function $L$
  slowly varying at $0$. Then $\mathrm{CV}_e^2\to 1/k$ as $e\to\infty$, so under linear
  cost $\varepsilon(e)\to 1/(1+k)$. If instead $G$ puts no mass near zero,
  $\underline\rho\equiv\operatorname{ess\,inf}\rho>0$, then $\mathrm{CV}_e\to 0$ and
  $\varepsilon(e)\to 0$.
\item[(d)] (Gamma case.) If $G$ is gamma with shape $k$, then $G_e$ is gamma with the
  same shape for every $e$, so $\mathrm{CV}_e^2=1/k$ and $\varepsilon=1/(1+k)$ at every
  effort level. This is Corollary~1(iii).
\end{enumerate}
\end{proposition}

Proposition~\ref{prop:cv} answers the question of which results depend on the
functional form. That offsetting is partial depends on nothing but heterogeneity in
how easy defects are to find. Its size is a single moment of a measurable object, the
detection rates of the defects that verification has not yet found. The constant
$1/(1+k)$ of Proposition~2 is the gamma case. More generally it is the high-effort
limit for any technology whose density of detection rates behaves like
$\rho^{k-1}$ near zero, so $k$ is best read as an index of how much mass sits on
defects that are almost impossible to find. Where testing eventually exhausts the
supply of hard defects, as when detection rates are bounded away from zero,
high-effort verification reverts to complete offsetting. The same happens, slowly,
when the density near zero vanishes faster than any power: for lognormal detection rates
with $\sigma=1$, $\varepsilon(e)$ falls from about $0.6$ at low effort to $0.15$ at
$e=1000$ (Online Appendix C).

% ---------------------------------------------------------------------
%  AI-specific extension: tool quality that changes the mix of errors
% ---------------------------------------------------------------------

\paragraph{When better tools make subtler errors.}
A better code generator need not merely produce fewer defects; it may produce
different ones. Write detection rates as $\rho=\theta\xi$, where $\xi\sim G$ has a
fixed distribution and $\theta$ scales how easy the tool's defects are to detect, and
let tool quality move both intensity and detectability, $\theta\propto\lambda^{\eta}$.
The parameter $\eta>0$ means that better tools ($\lambda\downarrow$) make defects that are harder to
find ($\theta\downarrow$). We call this case \emph{plausible-but-wrong} code.

\begin{proposition}[Verification-adjusted quality]
\label{prop:eta}
Under linear verification cost, holding detectability fixed the elasticity of
expected residual defects with respect to defect intensity is $\varepsilon$
of Proposition~\ref{prop:cv}, and holding intensity fixed the elasticity with
respect to detectability is $-(1-\varepsilon)$. Hence
\[
  R^*\;\propto\;\lambda^{\varepsilon}\,\theta^{-(1-\varepsilon)}
  \quad\text{locally, and}\quad
  \frac{d\ln R^*}{d\ln\lambda}=\varepsilon-\eta\,(1-\varepsilon).
\]
A better tool reduces residual defects if and only if
$\eta<\varepsilon/(1-\varepsilon)=\mathrm{CV}_{e^*}^2$. Under gamma-mixed detection
this threshold is $1/k$, and with $\hat k\approx 0.63$ it is about $1.6$.
\end{proposition}

The index $\lambda^{\varepsilon}\theta^{-(1-\varepsilon)}$ is a
\emph{verification-adjusted} measure of tool quality. Benchmark pass rates measure
$\lambda$ alone. What users receive depends on $\lambda$ and $\theta$ together, with
weights fixed by the heterogeneity of detectability. The same heterogeneity that makes
offsetting partial also sets how far a tool can shift its errors toward the hard-to-find
before its improvements stop reaching users.

% =====================================================================
%  Online Appendix B: proofs
% =====================================================================

\begin{proof}[Proof of Proposition~\ref{prop:cv}]
(a) For $e>0$ differentiation under the integral is justified because
$\rho^m e^{-\rho e}$ is bounded on $(0,\infty)$ for each $m$. Thus
$S'(e)=-\int\rho e^{-\rho e}dG$ and $S''(e)=\int\rho^2e^{-\rho e}dG$, and
$h(e)=-S'(e)/S(e)=\int\rho\,dG_e=\mu_e$. Then
\[
  h'(e)=-\frac{S''}{S}+\Bigl(\frac{S'}{S}\Bigr)^2
       =-\Bigl(\int\rho^2dG_e-\mu_e^2\Bigr)=-\sigma_e^2 .
\]
Because $G_e$ is absolutely continuous with respect to $G$ with a strictly positive
density, $G_e$ is degenerate if and only if $G$ is.

(b) By (a), $-h'/h=\sigma_e^2/\mu_e$, so $a=g+\sigma_e^2/\mu_e$ and
$a+h=g+\sigma_e^2/\mu_e+\mu_e$. Substitute into $\varepsilon=a/(a+h)$ from
Proposition~1. With $g=0$,
$\varepsilon=\sigma^2/(\sigma^2+\mu^2)=\mathrm{CV}^2/(1+\mathrm{CV}^2)$.
Monotonicity in $g$ follows from $\partial\varepsilon/\partial g
=\mu/(g+\sigma^2/\mu+\mu)^2>0$.

(c) Let $Y_e=e\rho$ under $G_e$. For $y>0$,
\[
  \Pr(Y_e\le y)=\frac{\int_0^{y/e}e^{-\rho e}dG(\rho)}{\int_0^\infty e^{-\rho e}dG(\rho)}
  =\frac{\int_0^{y}e^{-u}\,dG(u/e)}{\int_0^\infty e^{-u}\,dG(u/e)} .
\]
Divide numerator and denominator by $G(1/e)$. Regular variation gives
$G(u/e)/G(1/e)\to u^{k}$ for each $u>0$, and Potter's bounds
(Bingham, Goldie and Teugels 1987, Thm.~1.5.6) give an integrable dominating function
once the factor $e^{-u}$ is applied, uniformly for large $e$. Dominated convergence then yields
$\int_0^\infty u^m e^{-u}dG(u/e)/G(1/e)\to\int_0^\infty u^m e^{-u}\,d(u^k)
=k\,\Gamma(k+m)$ for $m=0,1,2$. (This is the Abelian half of Karamata's Tauberian
theorem; Feller 1971, \S XIII.5.) Hence $E[Y_e^m]\to\Gamma(k+m)/\Gamma(k)$, that is,
$Y_e$ has the moments of a Gamma$(k,1)$ variable in the limit. Since $\mathrm{CV}$ is
scale-invariant, $\mathrm{CV}_e^2=\mathrm{Var}(Y_e)/E[Y_e]^2\to
\bigl(k(k+1)-k^2\bigr)/k^2=1/k$.

If $\underline\rho>0$, fix $\delta>0$. On $(\underline\rho+\delta,\infty)$ the tilt
$e^{-\rho e}$ is at most $e^{-(\underline\rho+\delta)e}$, and on
$[\underline\rho,\underline\rho+\delta/2]$ it is at least $e^{-(\underline\rho+\delta/2)e}$, so
\[
  \frac{G_e\bigl((\underline\rho+\delta,\infty)\bigr)}
       {G_e\bigl([\underline\rho,\underline\rho+\delta/2]\bigr)}
  \le\frac{e^{-\delta e/2}}{G\bigl([\underline\rho,\underline\rho+\delta/2]\bigr)}\to0,
\]
where the denominator is positive by definition of the essential infimum. Hence $G_e$
concentrates on $\underline\rho$, so $\sigma_e\to0$ while $\mu_e\to\underline\rho>0$, and
$\mathrm{CV}_e\to0$.

(d) If $dG\propto\rho^{k-1}e^{-\beta\rho}$ then
$dG_e\propto\rho^{k-1}e^{-(\beta+e)\rho}$, a gamma density with shape $k$ and rate
$\beta+e$, so $\mathrm{CV}_e^2=1/k$ for all $e$.
\end{proof}

\begin{proof}[Proof of Proposition~\ref{prop:eta}]
Let $x=\theta e$ denote effort in units of the tool's detectability and
$S_\xi(x)=\int e^{-\xi x}dG(\xi)$, so $S(e)=S_\xi(\theta e)$. With $C(e)=ce$ the
developer chooses $x$ to minimise $(c/\theta)x+A\lambda s\,S_\xi(x)$, whose
first-order condition $c=A\lambda s\,\theta\,(-S_\xi'(x))$ depends on $(\lambda,\theta)$
only through $Z=\lambda\theta$. Hence $x^*=x^*(Z)$ and
$\ln R^*=\ln\lambda+\ln s+\ln S_\xi(x^*(Z))$. Let
$\kappa\equiv d\ln S_\xi(x^*)/d\ln Z$. Holding $\theta$ fixed,
$\partial\ln R^*/\partial\ln\lambda=1+\kappa$, which equals $\varepsilon$ by
Proposition~1. Holding $\lambda$ fixed,
$\partial\ln R^*/\partial\ln\theta=\kappa=\varepsilon-1$. With
$\theta\propto\lambda^\eta$, $d\ln R^*/d\ln\lambda=\varepsilon+\eta(\varepsilon-1)$,
which is positive if and only if $\eta<\varepsilon/(1-\varepsilon)$. By
Proposition~\ref{prop:cv}(b) the right side equals $\mathrm{CV}_{e^*}^2$, where the survivor
distribution is taken in units of $\xi$; the CV is scale-free.
\end{proof}

% ---------------------------------------------------------------------
%  Remark for Section 7: nonparametric estimation
% ---------------------------------------------------------------------

\paragraph{Remark (estimating $\varepsilon$ without a parametric form).}
For defect $j$ in a task with $M$ test inputs, of which $f_j$ fail, a random test
order leaves the defect undetected after a fraction $u$ of the suite with probability
approximately $(1-f_j/M)^{uM}=\exp(-\hat\rho_j u)$, where
$\hat\rho_j=-M\ln(1-f_j/M)$. (Exact hypergeometric survival can be used instead.)
The survivor CV at effort $u$ is the CV of $\{\hat\rho_j\}$ weighted by
$\exp(-\hat\rho_j u)$, so
\[
  \hat\varepsilon(u)=\frac{\widehat{\mathrm{CV}}_u^2}{1+\widehat{\mathrm{CV}}_u^2},
\]
with intervals from the task bootstrap already in use. Defects that fail every input
($f_j=M$) have $\hat\rho_j=\infty$ and weight zero for any $u>0$. We recommend reporting
the curve $\hat\varepsilon(u)$ as the primary estimate and the gamma fit as a summary.
Part (c) says what to expect: $\hat\varepsilon(u)$ should flatten at $1/(1+k)$ if the
lower tail of $\hat\rho_j$ is power-like, and drift toward zero if it is thin. The
lower-tail index $k$ can be estimated directly, for example with a Hill-type estimator on
the smallest values of $f_j/M$.

% References to add:
% Bingham NH, Goldie CM, Teugels JL (1987) Regular Variation (Cambridge University Press).
% Feller W (1971) An Introduction to Probability Theory and Its Applications, Vol. II, 2nd ed. (Wiley).
